\begin{lemma}[Part III, Algorithm (A) and invariant (6)]
Let \(T\) be an ETT with terminal PE index \(q\), terminal supporting
vertex \(v_n\), first-appearance slot \(j\), and a proper coloring
\(\sigma_n\) leaving exactly \(e_0\) uncolored.  Put
\(L^*=\noderef{ExtensionBaseRestorationIndices}\).  Assume the following
explicit conclusions of the preparations preceding Algorithm (A).
Here terminality means \(q+1=n\), the certified slot \(q.\mathsf{succ}\)
has kind PE and exact chosen vertex \(v_n\), and
\(j+1=m(v_n)\).
\begin{enumerate}
\item Every \(t\in L^*\) has exact certified data
      \((v_t,\gamma_t,\delta_t)\), one has \(v_n\prec v_t\), and \(T_j\)
      is an ordered initial segment of \(T_t\).  These are the order
      consequences proved by \noderef{ExtensionBaseSupportOrder}.
\item Every active nonterminal pair is \(T_t\)-strongly interchangeable
      under \(\sigma_n\), as in
      \noderef{ExtensionBaseClaimStrongInterchangeable}.
\item For the largest index \(t\) of each support class contained in
      \(L^*\), \(\delta_t\) is missing and \(\gamma_t\) is present at
      \(v_t\) under \(\sigma_n\).  This is the endpoint content of
      \noderef{ExtensionBaseClaimMissingColors}.
\item For \(s<t\) in \(L^*\), the exact pairs
      \(S_s=\{\gamma_s,\delta_s\}\) and
      \(S_t=\{\gamma_t,\delta_t\}\) intersect if and only if their
      supports agree and no member of that support class lies strictly
      between \(s\) and \(t\).  In that case
      \(S_s\cap S_t=\{\gamma_t\}=\{\delta_s\}\).  This is the exact
      overlap law of \noderef{ExtensionBaseStoredPairOverlap}.
\end{enumerate}
Then there is an \noderef{ExtensionBaseRestorationTrace} which enumerates
\(L^*\) in strictly decreasing order, starts at \(\sigma_n\), and
satisfies the \noderef{ExtensionBaseRestorationInvariant}.  In
particular, before every switch its chosen component is a path meeting
\(T_t\) exactly at \(v_t\), the color \(\delta_t\) is missing there, and
after the switch every still-unprocessed active nonterminal pair remains
strongly interchangeable.
\end{lemma}
\begin{proof}
Because \(L^*\) is finite, list its exact slot identities in strictly
decreasing order
\[
                         t_0>t_1>\cdots>t_{r-1}.               \tag{1}
\]
The exact-data assumption chooses \(v_{t_a},\gamma_{t_a},\delta_{t_a}\)
from the corresponding certified PE fields; no parallel edges are
identified in this choice.  Define \(\sigma^0=\sigma_n\), and recursively
define
\[
 \sigma^{a+1}=\noderef{TwoColorComponentSwap}
   (\sigma^a,\gamma_{t_a},\delta_{t_a},v_{t_a}).                \tag{2}
\]
Equations (1)--(2), together with the exact range equality, give all
fields of an \noderef{ExtensionBaseRestorationTrace}.  It remains to
prove the invariant by induction on \(a\).

Fix a support class \(L(w)\subseteq L^*\), and write its members in
increasing order as \(i_0<\cdots<i_p\).  The overlap assumption gives
\[
          \gamma_{i_{b+1}}=\delta_{i_b}\quad(0\leq b<p),       \tag{3}
\]
and says that all other intersections among its exact pairs are empty;
pairs belonging to different support classes are also disjoint.  At the
first processed member \(i_p\), the endpoint assumption says that
\(\delta_{i_p}\) is missing and \(\gamma_{i_p}\) is present at \(w\).
Switching their component therefore replaces the missing color
\(\delta_{i_p}\) at \(w\) by \(\gamma_{i_p}=\delta_{i_{p-1}}\).
Indeed, expand \noderef{MissingColors}: properness gives at most one
incident edge of either color, the missing delta gives none of color
\(\delta_{i_p}\), and the present gamma gives one of color
\(\gamma_{i_p}\); equation (2) exchanges the colors on exactly that
component.  Repeating this argument down the class and using (3) proves
that, immediately before processing \(i_b\), \(\delta_{i_b}\) is missing
and \(\gamma_{i_b}\) is present.  Disjointness of the other exact pairs
shows that switches in other classes do not alter either assertion.
This proves the missing-color part of the induction simultaneously for
all classes.

We next prove the precise component assertion at step \(a\).  By the
induction hypothesis, the pair at \(t_a\) is strongly interchangeable
under \(\sigma^a\), and the preceding paragraph says that its delta is
missing and its gamma is present at \(v_{t_a}\).  The finite component
has an edge and a vertex of incidence degree one, so the spanning path
from \noderef{TwoColorComponentSpanningPath} witnesses
\noderef{TwoColorPathComponent}.  If another vertex
\(x\in V(T_{t_a})\setminus\{v_{t_a}\}\) belonged to this component,
strong interchangeability would make the same connected component a
cycle.  A path component has an endpoint with subsingleton neighbor set,
whereas every vertex of the asserted cycle has two distinct cycle
neighbors, a contradiction.  Hence the component meets \(T_{t_a}\)
in precisely \(v_{t_a}\).  The terminal-order and initial-segment clauses
are the first explicit preparation and are independent of the coloring.

The component switch preserves properness by
\noderef{TwoColorComponentSwapProper}.  It merely interchanges two
\(\mathsf{some}\)-values and therefore preserves the statement that
exactly \(e_0\) is uncolored.  Now let \(s\ne q\) be an active slot not
processed through step \(a\), with exact support and colors.  If its
support differs from \(v_{t_a}\), the overlap law makes its pair disjoint
from the switched pair.  If its support agrees, decreasing order (1)
gives \(s<t_a\).  Thus the two alternatives in
\noderef{StrongInterchangeabilityKempePreserved} apply, and that lemma
shows that the pair at \(s\) remains strongly interchangeable under
\(\sigma^{a+1}\).  This is exactly the second half of invariant (6).

Initially the required interchangeability is assumption 2.  The
preceding paragraphs establish the induction step, including properness,
the exact-one-uncolored property, the missing endpoint color, the exact
tree intersection, and preservation for every unprocessed pair.  After
the last step the same induction statement gives the final clause; when
\(L^*=\varnothing\), the trace has length zero and it follows directly
from assumption 2.  Hence the constructed trace satisfies every clause
of \noderef{ExtensionBaseRestorationInvariant}.
\end{proof}
